Esta tesis desarrolla un sistema de inteligencia de negocio crediticio para instituciones de la economía social y solidaria del Ecuador. Se comparan tres modelos de Machine Learning para la prediccion multi-horizontede riesgo de incumplimiento. Los resultados muestran que LightGBM supera significativamente a las redes neuronales (AUC-ROC: 80\% vs 50\%), debido a las limitaciones de las redes neuronales con datos desbalanceados. El sistema usa de pipelines de MLOps automatizados con Apache Airflow, un datamart analítico en esquema estrella, registro de experimentos en MLFlow, y dashboards interactivos en Apache Superset para la toma de decisiones. La principal contribución es demostrar que los modelos basados en arboles de decisión son superiores para datos tabulares financieros con secuencias cortas.


\noindent\textbf{Palabras Clave}:

Riesgo crediticio; Deep Learning; CNN multi-horizonte; Datamart analítico; Apache Superset; Inteligencia Artificial; Business Intelligence; MLOps.
